\chapter{Synthesis: Integrating Formal Specifications into the Development Lifecycle}\label{ch:synthesis}

\section{The arc of the thesis}

The preceding five chapters were written, and in most cases submitted, as
self-contained studies. Read in isolation, each answers a bounded question:
whether inferred contracts help a language model write tests
(\cref{ch:article1}); whether such contracts can travel with the artefacts and
interfaces developers exchange (\cref{ch:article2}); whether a compositional
second pass recovers specifications a local pass misses (\cref{ch:article3});
whether inference can run unobtrusively in continuous integration
(\cref{ch:article4}); and what an end-to-end, verification-centric pipeline for
agentic development might look like (\cref{ch:article5}). This chapter is the
intellectual centre of the thesis. Its purpose is to show that the five studies
are not five papers stapled together but a single, cumulative argument, and to
answer the programme-level research questions with evidence marshalled across
chapters rather than within any one of them.

The thesis statement is this: \emph{formal specifications can be made cheap
enough, through automated inference, to serve as the always-present verification
spine of modern software development---including the AI-native, language-model-agent
lifecycle now emerging---and the same inferred contracts, when propagated
compositionally across a call graph, give a practical, sound (if incomplete)
account of behavioural version-compatibility between a library and its clients.}
Two concerns are entangled in that sentence by design. The first is the
classical adoption problem of formal methods: specifications are valuable but
expensive to write, so they are not written, so the value is never realised. The
second is the library/version-compatibility problem: a change to a dependency can
silently break a client whose behaviour depended on the old contract, and the
field has no cheap, behavioural account of when this happens. The thesis argues
that a single instrument---heuristic specification \emph{inference}, made
trustworthy by formal verification and made compositional by
weakest-precondition propagation---addresses both, because the act of inferring a
contract cheaply is also the act of producing the artefact whose change can be
propagated and checked.

The instrument is concrete throughout. Specifications are written in the Java
Modeling Language~\citep{leavens2006jml}, a behavioural-interface
specification language in the design-by-contract tradition~\citep{meyer1992dbc}.
They are inferred from source by a heuristic engine, the JML-Inferrer, that
reasons over the abstract syntax tree. They are verified by OpenJML's extended
static checker~\citep{cok2014openjml}, which discharges proof obligations through
the Z3 satisfiability-modulo-theories solver~\citep{demoura2008z3}. Every claim
about whether an inferred specification is \emph{correct} is grounded in this
verification path rather than in inspection, and this matters for the integrity
of the whole programme: a heuristic that guesses a contract is only as useful as
the mechanism that can certify the guess.

\subsection{How the five studies compose}

The studies compose as a single pipeline of claims, each chapter discharging a
precondition for the next.

\Cref{ch:article1} establishes that inference is \emph{worth doing}. It is the
foundational premise: that automatically inferred specifications, although
heuristic and partial, are good enough to be useful downstream. Supplying
inferred JML contracts to a language model improves the unit tests it generates,
converting the abstract value of specifications into a measured, downstream
benefit. Without this result, the rest of the programme would be an exercise in
producing artefacts of unproven worth.

\Cref{ch:article2} establishes that inference is \emph{deployable}. A
specification that lives in a separate document goes stale; to become part of a
lifecycle it must travel with the code. This chapter shows that JML contracts can
be embedded into compiled Java artefacts and into OpenAPI interface documents and
recovered again at full round-trip fidelity, giving inferred contracts a carriage
mechanism through the very artefacts teams already exchange.

\Cref{ch:article3} establishes that inference is \emph{extensible and
compositional}. Single-pass heuristic inference is local by construction: it
reasons about a method in isolation. A second, compositional pass that propagates
weakest-precondition obligations across method boundaries strictly extends the
local pass, recovering a substantial number of additional well-formed
preconditions on a large corpus. This result is load-bearing twice over. It
strengthens the inference instrument, and it supplies the mechanism on which
version-compatibility analysis later rests, because the reasoning that recovers a
caller's obligation from a callee's contract is exactly the reasoning that
propagates a contract \emph{change} across a call graph.

\Cref{ch:article4} establishes that inference is \emph{a process, not a tool}. It
integrates the pipeline into continuous integration---GitHub Actions, GitLab CI,
and Jenkins---under three principles: per-pull-request, diff-only analysis with
content-hash caching; fail-open semantics so inference never blocks a merge; and
persistent, in-place reporting. This is the step from an executable into a
routine, low-friction part of the development cycle.

\Cref{ch:article5} is the capstone. It draws the four threads---cheap, deployable,
compositional, routine---into a single verification-cen\-tric reference pipeline
for the lan\-guage-model-agent lifecycle, in which a machine-check\-able contract
is the pivot between ambiguous human intent and precise machine verification,
serving simultaneously as the input an agent builds against and the oracle its
output is checked against.

The shape of the argument, then, is: inference made cheap (\cref{ch:article1})
$\rightarrow$ shown useful and deployable (\cref{ch:article1,ch:article2})
$\rightarrow$ carried with artefacts and run in continuous integration
(\cref{ch:article2,ch:article4}) $\rightarrow$ compositionally refined and
propagated (\cref{ch:article3}) $\rightarrow$ integrated as the verification spine
of the lifecycle (\cref{ch:article5}). The remainder of this chapter answers the
research questions against this arc, presents the integrated reference lifecycle
as the synthesis, draws the explicit link between compositional propagation and
behavioural version-compatibility, articulates the development-process
contribution, and closes with a programme-wide treatment of validity and
limitations.

It is worth dwelling briefly on why this particular ordering is the right one,
because the dependency between the chapters is logical and not merely
chronological. A reader sympathetic to formal methods might ask why the programme
did not begin with the compositional analysis, which is its most technically
ambitious component. The answer is that compositional propagation is only
meaningful if the local contracts it propagates are themselves worth having, and
the worth of an \emph{inferred} contract---as opposed to a hand-written one---is
exactly what \cref{ch:article1} had to establish first. A sceptic of the whole
enterprise is entitled to doubt that a heuristic, machine-generated contract
carries enough information to be useful at all; until that doubt is answered,
every later result rests on sand. Equally, the carriage and continuous-integration
results (\cref{ch:article2,ch:article4}) are not optional engineering footnotes but
necessary conditions for the central thesis claim: a contract that cannot survive
the journey from where it is inferred to where it is needed, or that cannot be
produced without a developer stopping to invoke a tool, will never become the
\emph{always-present} spine the thesis argues for. The phrase ``always present''
does real work in the thesis statement, and it is the carriage and process chapters
that earn it. Only once a contract is cheap, useful, durable, and routinely
produced does it make sense to ask, as the capstone does, what an entire lifecycle
built around such contracts would look like.

\section{Answering the research questions}

The programme is organised around five research questions. Each is answered below
with evidence drawn from the relevant chapters, and each answer is qualified
honestly by the \emph{kind} of evidence that supports it: a measurement, a
feasibility demonstration, or a design validated by instantiation rather than by a
controlled effect size.

\subsection{RQ1: inference fidelity and compositional reach}

\emph{How faithfully can specifications be inferred from unannotated code, and how
far can a compositional pass extend that reach?}

The fidelity claim is anchored to formal verification rather than to inspection.
Across the corpora used in \cref{ch:article1,ch:article3}, a meaningful fraction
of inferred contracts are not merely plausible but \emph{machine-checked}: OpenJML
discharges them against the implementation via extended static checking. This is a
stronger guarantee than the inductive-invariant detection of dynamic tools such as
Daikon~\citep{ernst2007daikon}, whose candidate invariants are likelihoods over
observed runs rather than proved properties. The fidelity is partial---the engine
is heuristic and emits no contract where it cannot recognise a pattern---but
where it emits one and verification accepts it, the contract is sound with respect
to the code.

The compositional-reach claim is the central quantitative result of
\cref{ch:article3}: a compositional pass that propagates weakest-precondition
obligations across call boundaries strictly extends single-pass inference,
yielding a substantial number of additional well-formed preconditions on a large
corpus that the local pass could not recover. Crucially, the figures reported are
the \emph{well-formedness-filtered} figures: preconditions are admitted only when
their entry-state scope is satisfiable, which is why the headline count is the
conservative one rather than the raw, unfiltered total.

A second dimension of fidelity deserves emphasis because it distinguishes this
work from the broad family of specification-mining techniques. Dynamic invariant
detection observes a program running over a test suite and reports the properties
that held on every observed run; the resulting invariants are \emph{empirical
generalisations}, true of the executions seen and conjectured of the rest. The
JML-Inferrer's contracts, by contrast, are subjected to extended static checking,
which reasons over \emph{all} executions symbolically. The two approaches are
complementary rather than competing---a dynamic tool will propose richer
candidate properties, while a static checker certifies a narrower set
unconditionally---but the thesis deliberately stakes its fidelity claim on the
stronger, exhaustive guarantee, because the downstream uses (an oracle that judges
generated code, a contract that bounds a version-compatibility analysis) require
properties that are true rather than merely probable. The cost of that choice is
silence: the engine emits nothing where it cannot find a pattern the checker will
accept, and so RQ1's fidelity is necessarily reported as a rate over emitted
contracts, not as a coverage of all behaviours.

Strength of evidence: this RQ is answered by \emph{measurement}. The fidelity and
reach numbers are counts and rates over real corpora, verified by an independent
checker. The principal caveat is corpus dependence, addressed under threats to
validity below.

\subsection{RQ2: downstream utility for test generation}

\emph{Do inferred specifications make a language model a better test author?}

\Cref{ch:article1} answers in the affirmative through a controlled comparison:
supplying inferred JML contracts to a language model improves the unit tests it
generates relative to the same model without them. The mechanism is the
test-oracle problem~\citep{barr2015oracle}---a generated test exercises code but
needs an independent statement of intended behaviour to judge the result---and an
inferred contract supplies exactly that statement. This connects to a live concern
in language-model code generation: models hallucinate
behaviour~\citep{zhang2025hallucination} and generate plausible but unfounded
oracles~\citep{liu2026oracle}, and benchmarks such as
EvalPlus~\citep{liu2023evalplus} show that surface-passing generated code often
fails under strengthened tests. A contract gives the model a fixed point against
which to anchor.

Strength of evidence: this RQ is answered by a \emph{controlled measurement}---a
direct comparison with and without the treatment---but the effect is demonstrated
on a bounded set of programs and one family of models, so it establishes
direction and plausibility rather than an industry-wide effect size.

\subsection{RQ3: carriage and process integration}

\emph{Can inferred specifications be carried with the artefacts developers exchange
and integrated into the development process without disrupting it?}

This question has two halves, answered by two chapters. The carriage half
(\cref{ch:article2}) demonstrates full round-trip embedding of JML into compiled
Java artefacts and into OpenAPI interface documents: a contract inferred at one
point in the lifecycle survives compilation, distribution, and recovery without
loss. The process half (\cref{ch:article4}) demonstrates that inference runs as a
first-class continuous-integration step under fail-open, diff-only, cached,
in-place-reporting discipline, so that specification recovery becomes a routine
event on every pull request rather than a special activity.

Strength of evidence: this RQ is answered by \emph{feasibility demonstration with
fidelity measurement}. The round-trip fidelity is measured exactly (specifications
recovered are identical to those embedded); the CI integration is demonstrated to
work across three platforms. What is \emph{not} claimed is a measured improvement
in team outcomes, which would require a field study and is left as future work.

\subsection{RQ4: compositional behavioural version-compatibility}

\emph{Can compositional propagation of inferred contracts give a behavioural
account of when a library change breaks a client?}

This is the half of the thesis core concerned with library/version compatibility,
and it is answered by combining the mechanism of \cref{ch:article3} with the
propagation analysis developed in \cref{ch:article5}. The argument is: if a
callee's contract can be inferred, and the weakest-precondition obligation it
imposes on a caller can be computed compositionally, then a \emph{change} to the
callee's contract has a computable \emph{behavioural blast-radius}---the set of
callers whose obligations the change would now violate. A dependency upgrade
becomes a contract-diff event whose effect on clients can be checked rather than
discovered in production. This is a behavioural notion of compatibility, distinct
from and finer-grained than the syntactic, library-maintainer-declared semantics
of version numbers~\citep{preston2013semver}; it sits alongside empirical accounts
of how breaking changes actually arise~\citep{ochoa2022breaking,mostafa2017behavioral}
and how developers reason about them~\citep{jayasuriya2024understanding}.

Strength of evidence: this RQ is answered by a \emph{design validated by
instantiation}, with an important honesty qualification. The propagation is
\emph{sound in flagging breaks}: when it reports that a contract change violates a
caller's obligation, the violation is real with respect to the inferred contracts.
It is \emph{not complete}: it will miss breaks when a third-party contract is
inferred too weakly to capture the depended-upon behaviour, or when a Z3 timeout
prevents discharge of the relevant obligation. A missed break is therefore always
possible; the analysis raises an alarm reliably but cannot guarantee silence means
safety.

\subsection{RQ5: lifecycle integration}

\emph{Can inferred, verified, compositional specifications serve as the contract
layer and verification loop of an AI-native development lifecycle?}

This is the synthesising question, answered by \cref{ch:article5} drawing on all
that precedes it. The capstone proposes a reference pipeline in which a formal
contract is the dual-role pivot of the lifecycle---input and oracle at once---and
identifies the three things that current agentic, specification-driven lifecycles
lack: a machine-checkable contract (rather than prose intent), a closed
counterexample-guided verification loop, and an independent authority for the
oracle so that the agent does not grade its own work. The pipeline supplies all
three, and reuses the validated instruments of the earlier chapters as its trusted
components: inference for the contract, embedding for its carriage, compositional
propagation for change analysis, and continuous integration as its execution
substrate.

Strength of evidence: this RQ is answered by a \emph{design validated by
instantiation rather than by a controlled trial}. The architecture is concrete and
partially implemented, and its components are individually evidenced by the earlier
chapters; the end-to-end lifecycle is demonstrated for feasibility and realism, not
established by a controlled effect size. The industrial pilot that would supply
such an effect size is explicitly future work.

\section{The integrated reference lifecycle as the synthesis}

The synthesis is best stated as a single picture of how the parts fit together
once they are composed. The organising idea is the \emph{dual-role
specification}: one machine-checkable formal contract that plays two roles at once.
As \emph{input}, it is the precise statement of intent against which a generative
agent (or a human) writes an implementation. As \emph{oracle}, it is the property
against which that implementation is verified. The contract is therefore the pivot
between two registers that have never sat comfortably together: the ambiguous,
natural-language intent of a human request, and the precise, decidable obligation a
verifier can discharge. Specification-driven agent toolkits already gesture at this
pivot---GitHub's Spec Kit~\citep{githubspeckit}, AWS's specification-driven
agentic development~\citep{awsaidlc}, and the Kiro environment~\citep{awskiro} all place a
specification at the centre---but they keep the specification in prose, which a
machine cannot check, so the loop never closes~\citep{ozkaya2023llmse}.

Three further properties make the lifecycle work, and each is inherited from an
earlier chapter. \emph{Inference is the bootstrap.} The dual-role contract is only
viable if it is cheap to obtain; the authoring cost that has always blocked
design-by-contract adoption is removed by inferring the contract from existing code
(\cref{ch:article1,ch:article3}) rather than asking a developer to write it. This
is the adoption half of the thesis core, discharged. \emph{Carriage makes the
contract durable.} The contract must survive the round-trip through build and
distribution to be available at every lifecycle stage (\cref{ch:article2}).
\emph{Assurance is tiered.} Not every contract is verified to the same depth: some
are machine-checked end to end, some are runtime-checkable only, and some are
heuristic suggestions awaiting ratification. The lifecycle is explicit about which
tier a given contract occupies, so that trust is calibrated rather than uniform.

\subsection{The independent-authority constraint and its honest limit}

The most delicate property is the \emph{independent authority} of the oracle. The
hazard in any self-improving generative loop is circularity: if the same model
both writes the code and judges it, the judgement inherits the author's blind
spots, and the system can certify its own mistakes. Recent work on closing the
language-model verification loop confronts exactly this---Clover couples code,
documentation, and tests into a consistency check~\citep{sun2024clover}, and oracle
soundness for generated assertions is an active concern~\citep{liu2026oracle}. The
reference pipeline addresses circularity by deriving the oracle from an authority
independent of the code-writing agent: a contract inferred and \emph{verified} by
the formal toolchain, not asserted by the generator.

It is important to state precisely what this does and does not buy. The
independent-authority constraint removes \emph{self-grading and authorship
circularity}: the entity that proves the code correct is not the entity that wrote
the code, and the proof is a re-checkable artefact rather than the author's
assurance. What it does \emph{not} remove is \emph{shared-misreading correlation}.
If a single model is used both to help recover the contract and to write the code,
and that model misreads the underlying intent in the same way in both roles, the
contract and the code can be consistently wrong together, and an independent
verifier will happily certify their agreement. The pipeline's answer to this
residual risk is not a technical one but a procedural one: \emph{human
ratification} of the contract at the tier boundary, so that a person confirms the
contract states the intended behaviour before it is promoted to an authoritative
oracle. The thesis is candid that this is a mitigation of a correlated-error risk,
not its elimination.

\Cref{fig:synthesis-lifecycle} summarises the lifecycle at a high level.

\begin{figure}[t]
\centering
\begin{tikzpicture}[
  node distance=9mm and 14mm,
  box/.style={draw, rounded corners, align=center, inner sep=4pt,
              minimum height=8mm, font=\small},
  contract/.style={box, draw=black, very thick, fill=black!6},
  arr/.style={-{Stealth[length=2.2mm]}, thick},
]
  \node[box] (intent) {Human intent\\(ambiguous,\\natural language)};
  \node[contract, right=of intent] (spec) {Formal contract\\(dual role:\\input \& oracle)};
  \node[box, right=of spec] (code) {Implementation\\(agent / human)};
  \node[box, below=14mm of code] (verify) {Verifier\\(OpenJML + Z3)};
  \node[box, below=14mm of spec] (infer) {Inference\\(bootstrap)};

  \draw[arr] (intent) -- (spec);
  \draw[arr] (spec) -- node[above,font=\scriptsize]{input} (code);
  \draw[arr] (code) -- (verify);
  \draw[arr] (spec.south) |- (verify.east) node[pos=0.75,below,font=\scriptsize]{oracle};
  \draw[arr] (infer) -- (spec);
  \draw[arr] (infer.west) -| (intent.south);
  % counterexample-guided loop
  \draw[arr] (verify.west) -| ($(code.south)+(0,-9mm)$) -- ++(0,0)
        coordinate (cx) ;
  \draw[arr] (verify) -- ++(-3.2,0) |- (code.south)
        node[pos=0.25,left,font=\scriptsize]{counterexample};
  \begin{scope}[on background layer]
    \node[draw, dashed, rounded corners, fit=(spec)(code)(verify)(infer),
          inner sep=8pt, label={[font=\scriptsize]above:verification-centric loop}] {};
  \end{scope}
\end{tikzpicture}
\caption[The integrated reference lifecycle]{The integrated reference lifecycle.
Inference bootstraps a machine-checkable formal contract from existing code and
ambiguous intent; the contract then plays a dual role---input to the implementing
agent and oracle for the verifier---and the verifier returns counterexamples that
close the loop. Independent authority (the verifier is not the code author) and
tiered, human-ratified assurance govern when a contract is trusted.}
\label{fig:synthesis-lifecycle}
\end{figure}

\Cref{fig:synthesis-lifecycle} makes the closed loop visible: the same contract
node feeds both the implementer and the verifier, the verifier returns
counterexamples to the implementer, and inference is the off-ramp that seeds the
contract in the first place. This is the structural difference from prose-centred
agentic toolkits, in which there is an analogous box but no edge from a verifier
back to the implementer, because prose intent admits no verifier.

\section{Compositional propagation as behavioural version-compatibility}

The two halves of the thesis core---adoption and version-compatibility---meet at a
single mechanism, and it is worth stating the link explicitly because it is the
claim that makes the thesis one argument rather than two.

The mechanism is the compositional, weakest-precondition propagation of
\cref{ch:article3}. In that chapter it is used \emph{statically}, within a single
version of a codebase, to recover a caller's precondition from a callee's contract.
\Cref{ch:article5} reuses the identical mechanism \emph{differentially}, across two
versions of a dependency: given the old and new inferred contracts of a callee, the
propagation computes which callers' obligations the change disturbs---the
behavioural blast-radius. The same arrow that, run forwards within a version,
strengthens a caller's inferred specification, when run across versions becomes the
edge that carries a breaking change from a library to its clients. This is the
concrete bridge: compositional refinement is not merely \emph{analogous} to
compatibility analysis; it is the \emph{same computation} applied to a contract
delta.

A worked intuition makes the bridge concrete. Suppose a utility method
\texttt{clamp(int x, int lo, int hi)} is inferred, in version one of a library, to
carry the precondition \texttt{lo <= hi} and the postcondition that the result lies
in $[\mathtt{lo},\mathtt{hi}]$. A client that calls \texttt{clamp} and then indexes
an array with the result has, through compositional propagation, an inferred
obligation that the array is at least \texttt{hi+1} long---an obligation the local
pass over the client alone could never have produced, because it depends on the
callee's postcondition. Now suppose version two of the library relaxes
\texttt{clamp} so that it may return \texttt{lo} when \texttt{lo > hi} rather than
requiring the precondition. Run differentially, the same propagation machinery
observes that the client's array-bounds obligation, which was discharged under the
old contract, is no longer guaranteed under the new one, and flags the client as
within the blast-radius. No new analysis was invented for the differential case;
the version-one and version-two contracts were fed to the identical
weakest-precondition propagation, and the disturbance fell out of the comparison.
This is what it means to say compatibility analysis \emph{is} compositional
refinement applied to a contract delta.

The honesty of the claim turns on soundness versus completeness, and the thesis is
deliberate about which it has. The propagation is \emph{sound in flagging breaks}.
When it reports that a client's obligation is violated by a dependency change, the
violation is a genuine logical consequence of the inferred contracts; there are no
false alarms relative to those contracts. The propagation is \emph{not complete}.
It will fail to flag a real break in two well-characterised circumstances. First,
when the third-party callee's contract is inferred too weakly---the heuristic
recovered a vacuous or partial contract that does not capture the behaviour the
client actually depended on---the change to that behaviour is invisible because the
contract never expressed it. Second, when discharging the relevant obligation
exceeds the solver's reach---a Z3 timeout, a multiplicative or array-theory query
the solver cannot decide within budget---the obligation is left unproved and the
potential break unflagged. The practical reading is therefore asymmetric and must
be stated as such to users: an alarm is trustworthy, but silence is not a
guarantee of compatibility. This is the correct and defensible posture for a tool
built on heuristic inference and an undecidable underlying logic, and it is
preferable to a tool that claims completeness it cannot deliver.

\section{The development-process contribution}

Beyond the instruments, the thesis advances a claim about how software
\emph{development} itself changes when verified, inferred contracts are always
present. The contribution is a shift in the unit of review and the nature of
trust.

\paragraph{From code review to specification review.} Today, a change is reviewed
as a diff of code, and a reviewer reconstructs intent from the implementation. When
every change carries an inferred, verified contract, the reviewable unit becomes
the \emph{specification delta plus its proof}: the reviewer reads what the
behaviour is now required to be, and the verifier certifies that the code meets it.
Review shifts from inferring intent out of code to confirming intent stated as a
contract. This does not abolish code review---a reviewer may still object to
\emph{how} something is done---but it relocates the question of \emph{what} the code
must do into a checkable artefact.

\paragraph{Transferable trust.} The proof that a contract holds is a re-checkable
certificate, not the reviewer's eyeballing. Trust becomes transferable: a
downstream consumer of a library need not re-read its source to believe a property;
they can re-run the check. This is the property that makes the version-compatibility
half of the thesis operationally meaningful---a client trusts a dependency's
contract because the contract is certified, and a dependency upgrade becomes a
\emph{spec-diff event} whose blast-radius is computed, not a leap of faith mediated
by a version number.

\paragraph{A process-mining analogue.} The shape of this change has a precedent in
the business-process world, and naming it sharpens the contribution. Where process
mining discovers a model from event logs and conformance
checking~\citep{rozinat2008conformance,vanderaalst2016processmining,dumas2018bpm}
measures the fit between recorded behaviour and that model, the present programme
\emph{infers} a behavioural contract from code and then \emph{checks conformance}
of the implementation---and of its callers, across versions---against it. Inference
is the discovery step; OpenJML verification and compositional propagation are the
conformance-checking step. The analogy is not decorative: it places this work in a
recognised lineage of \emph{discover-then-conform} methods and clarifies that the
novelty is the medium (formal contracts over source code, machine-checked) rather
than the overall epistemic move.

\paragraph{An organisational note.} Insofar as it is relevant, the
version-compatibility result touches Conway's observation that systems mirror the
communication structures of the organisations that build them~\citep{conway1968committees}:
a contract-diff at a module boundary is, in effect, a coordination event between
the teams on either side of that boundary, and making the blast-radius computable
turns an implicit, cross-team coordination cost into an explicit, checkable one.
The thesis does not pursue the organisational dimension empirically---that would be
a different study---but notes that the mechanism it builds operates precisely at
the boundaries where Conway's law predicts coordination friction. Whether such a
mechanism is actually adopted is, in turn, a question of diffusion of
innovations~\citep{rogers2003diffusion}: the design choices of fail-open CI
integration (\cref{ch:article4}) and zero-authoring-cost inference are deliberate
reductions of the adoption barriers that diffusion theory identifies, namely
relative disadvantage and complexity.

\section{Threats to validity across the programme}

Each constituent study states its own threats; this section addresses those that
operate at the \emph{programme} level, where a weakness in one study could
compromise the cumulative argument. They are grouped by the standard construct,
internal, external, and conclusion categories~\citep{runeson2009casestudy,yin2018case},
and each is paired with its mitigation.

\subsection{Construct validity}

The central construct risk is that \emph{``a verified inferred contract'' may not
measure ``correct behaviour''}. A contract can be verified yet trivially weak: an
inference that emits \texttt{ensures true} is vacuously verifiable and says nothing.
The programme mitigates this in two ways. First, the well-formedness filtering of
\cref{ch:article3} excludes contracts whose entry-state scope is unsatisfiable, and
the headline reach figures are the filtered ones, so vacuous contracts are not
counted as wins. Second, the downstream-utility study (\cref{ch:article1}) measures
contracts by an \emph{external} effect---improved generated tests---which a vacuous
contract cannot produce, providing a check on construct validity that does not
depend on the contract's internal form. A residual risk remains: a contract can be
non-vacuous yet still under-specify intent, which is precisely why human
ratification sits at the tier boundary of the lifecycle.

\subsection{Internal validity}

Two internal threats dominate. The first is \emph{heuristic inference skew}: the
engine recognises patterns it was built to recognise and is silent elsewhere, so
the population of inferred contracts is biased toward the tractable. This could
inflate apparent fidelity if the easy cases are also the verifiable ones. The
mitigation is that fidelity is reported as a rate \emph{conditioned on emission and
verification}, and the silence of the engine is reported honestly rather than
counted as success; the thesis does not claim coverage it does not have. The second
is \emph{language-model non-determinism}: the test-generation result of
\cref{ch:article1} and any agentic-loop demonstration depend on a stochastic model
whose output varies run to run. The mitigation is the controlled, paired design of
the comparison---the same prompts with and without the contract treatment---so that
non-determinism is shared across conditions and the treatment effect is isolated;
the limitation is that a single model family was used, so the result is a direction,
not a universal constant.

\subsection{External validity}

The principal external threat is \emph{single-corpus generalisability}. Much of the
quantitative evaluation, especially the compositional-reach measurement, leans on
Apache Commons Lang and similarly structured Java libraries. A library written in a
markedly different idiom---heavy reflection, dynamic proxies, concurrency-dominated
control flow---may expose patterns the heuristics do not recognise and contracts the
solver cannot discharge. The mitigation is partial: the studies use real,
widely-depended-upon code rather than synthetic benchmarks, so the result is at
least representative of an important class of production Java; but the thesis does
\emph{not} claim that the reach figures transfer unchanged to arbitrary corpora, and
it identifies multi-corpus replication as the obvious next step. A second external
threat is that the AI-native lifecycle of \cref{ch:article5} targets a fast-moving
ecosystem; the specific agent toolkits it positions against~\citep{githubspeckit,awskiro,awsaidlc}
may evolve. The mitigation is that the architectural claims---dual-role contract,
closed loop, independent authority---are stated at a level above any particular
toolkit.

\subsection{Conclusion validity}

The dominant conclusion threat is \emph{solver decidability and scalability}. The
verification spine rests on Z3, and the underlying logic is undecidable; the
programme repeatedly encounters timeouts on multiplicative arithmetic, array-theory
queries, and rich preconditions. This bounds the generalisability of every claim
built on verification: where the solver cannot decide, the pipeline neither
confirms nor refutes. The mitigation is twofold and honest. Operationally, the
configuration is tuned (bigint specification arithmetic, safe code arithmetic, a
chosen solver version) to maximise the fraction discharged within budget, and the
fail-open and soundness-not-completeness postures ensure that an undecided
obligation degrades to ``unflagged/unverified'' rather than to a false guarantee.
Epistemically, the thesis states plainly that \emph{characterisation is not a
correctness proof}: where a result is a measurement of how the heuristics and
solver behave on a corpus, it is reported as such and not elevated to a theorem
about all programs.

\subsection{The capstone's deferred validation}

A final, cross-cutting threat is specific to \cref{ch:article5}: the integrated
lifecycle is \emph{partially built, with end-to-end validation deferred}. Its
components are individually evidenced by the earlier chapters, but the composed
system has not been subjected to a controlled trial. The mitigation is to claim
exactly what is demonstrated---feasibility and architectural coherence, with each
component validated in isolation---and to label the controlled lifecycle evaluation
as future work rather than as a completed result. The thesis would be weaker, not
stronger, if it overstated this.

There is also a subtler composition risk worth naming: the validity of a composed
system is not the conjunction of the validities of its parts. Each chapter's result
holds under that chapter's assumptions, and those assumptions are not identical. The
test-generation study assumes a particular interaction protocol with a language
model; the carriage study assumes a build toolchain that preserves embedded
metadata; the compositional study assumes that callee contracts are available and
non-vacuous; the continuous-integration study assumes a fail-open policy that some
organisations may reject. When these are stacked into one lifecycle, an assumption
satisfied in isolation may be violated in combination---for instance, a contract
strong enough to drive a verifier may be too verbose to embed efficiently, or a
fail-open CI policy may admit a weak contract that then under-constrains a
downstream blast-radius analysis. The thesis does not resolve these interaction
effects empirically; it identifies them as exactly the questions a future
end-to-end evaluation must measure, and treats the per-component evidence as
necessary but not sufficient for a claim about the whole. Naming the gap is itself
part of the contribution: it tells a future evaluator where to look.

\section{What the thesis does and does not claim}

It is worth closing the synthesis by drawing the boundary of the contribution
explicitly, because the value of a candid thesis is partly in its restraint.

The thesis \emph{claims}: that heuristic inference can produce formal contracts
cheaply enough to remove the authoring barrier to design-by-contract adoption; that
such contracts, although partial, verifiably hold against their code and are useful
downstream; that they can be carried with artefacts and run routinely in continuous
integration; that compositional propagation of inferred contracts gives a sound (if
incomplete) behavioural account of library/version compatibility; and that these
parts compose into a coherent, verification-centric reference lifecycle for
AI-native development in which a dual-role formal contract is the pivot between
human intent and machine verification.

The thesis \emph{does not claim} a controlled, industry-scale effect size for any
of this. The test-generation result is a controlled comparison on a bounded
program set and one model family, establishing direction and plausibility rather
than a population parameter. The carriage and continuous-integration results are
feasibility-and-fidelity demonstrations, not measurements of team productivity. The
version-compatibility analysis is sound in alarm but admittedly incomplete in
silence, bounded by weak third-party contracts and solver timeouts. The capstone
lifecycle is a design validated by instantiation, with its end-to-end controlled
evaluation---the industrial pilot---explicitly future work. And throughout,
characterisation of how the instruments behave on real corpora is reported as
characterisation, not as a proof of correctness over all programs.

What unifies these claims and disclaimers is the thesis's single move: making
formal specifications cheap by inferring them, trustworthy by verifying them, and
useful by carrying and propagating them through the lifecycle---so that the same
contract that lowers the barrier to formal methods also becomes the instrument for
reasoning about how software changes break the programs that depend on it. That
move is the contribution, and the five studies are its cumulative evidence.
